\section{Parameter weights and the criterion for recovering the algebra}
\label{sec:weights-and-parameters}

We continue under the assumption $\charac k=2$, with $A_0=kP$ ($P=\ray A_0$)
and $\C=\gr_\JJ\ind(\modu A)\simeq\ind(\modu A_0)$. In this section we compute
the universal group of additive weights on the ray category $P$, show how
ray weights extend to mesh weights on $\C$, and use the first-order module
lifting of Section~\ref{sec:positive-deformations} to embed the
positive-degree Hochschild cohomology $\HH^2(A_0)_{>0}$ into the mesh
cokernel $\bigoplus_{m\geq2}\operatorname{coker}D_m(\C)$. Crucially, the
weight lattice $L=W(P)$ separates the residual cohomology classes attached to
distinct periodic components even when they share the same categorical
degree.

\subsection{Normalizing the scalar mesh coefficients}

By Lemma~\ref{lem:unit-mesh-coefficients}, we may assume from now on that all
scalar coefficients in the mesh relations of $\C$ are equal to $1$:
explicitly, rescaling each arrow $a$ of $\Gamma_A$ by $\mu_a\in k^*$ and
dividing the mesh relation at each nonprojective vertex $v$ by $\eta_v\in k^*$
normalizes the coefficients $c_{v,a}$ to $1$ whenever
\begin{equation}
 \mu_{\sigma a}\mu_a=\eta_v/c_{v,a}\qquad\text{for every arrow }a:u\to v.
 \label{eq:global-mesh-scales}
\end{equation}

\subsection{The integral group of ray weights}

We now establish an integral counterpart of the scalar weight trivialization
of Lemma~\ref{lem:scalar-object-values}. This result holds in arbitrary
characteristic (compare the fundamental-group and coefficient analysis in
Bretscher and Gabriel \cite[Sections~2.11--2.12]{BretscherGabriel}) and will
also be used in Part~\ref{part:characteristic-two} to lift modulo-two
weights on smaller ray tables.

\begin{definition}
\label{def:scalar-weight-group}
For a finite ray category $Q$, its \emph{scalar-weight group} is the abelian
group
\[
 W(Q)=\Z\langle[r]:r\ne0,1\rangle/
       \langle[rs]-[r]-[s]:rs\ne0\rangle,
\]
where the classes of identity morphisms are set to zero. Thus homomorphisms
$W(Q)\to M$ into an abelian group $M$ correspond bijectively to additive
weights on the nonzero morphisms of $Q$ (and, for $M=k^*$ written
multiplicatively, $\Hom(W(Q),k^*)$ is the group of multiplicative scalar
weights $a(rs)=a(r)a(s)$ for $rs\ne0$).
\end{definition}

\begin{proposition}
\label{prop:integral-ray-weights}
Let $A_0=kP$ ($P=\ray A_0$) for a connected basic representation-finite
algebra $A$ with $n$ primitive idempotents, let $G=\pi_1(K)$ be the free deck
group of Subsection~\ref{subsec:universal-cover}, and let $\mathcal L$ be the
category of lifted projectives. Let
$I_{\mathcal L}\subset\Z^{(\operatorname{ob}\mathcal L)}$ denote the subgroup
of formal integer combinations of objects of $\mathcal L$ whose coefficients
sum to zero.
\begin{enumerate}[label={\upshape(\arabic*)}]
\item Sending a lifted ray $r\in\mathcal L(x,y)$ to the difference
$e_y-e_x\in I_{\mathcal L}$ induces an isomorphism
$W(P)\cong(I_{\mathcal L})_G$ with the group of $G$-coinvariants of
$I_{\mathcal L}$. Under this identification, the coinvariant class
$\theta_g=[e_{gx}-e_x]$ is independent of $x\in\operatorname{ob}\mathcal L$,
and the map $g\mapsto\theta_g$ embeds the abelianization $G_{\mathrm{ab}}$
as a direct summand of $W(P)$ with free complement $\Z^{n-1}$:
\[
                       W(P)\cong\Z^{n-1}\oplus G_{\mathrm{ab}}.
\]
In particular, since $G=\pi_1(K)$ is finitely generated, $W(P)$ is a free
abelian group of finite rank.
\item Every additive weight on $\mathcal L$ with values in an abelian group
is a difference of object values.
\end{enumerate}
\end{proposition}
\begin{proof}
For any finite directed thin representation-finite category $T$, let $U(T)$
denote the abelian group generated by symbols $[r]$ for the nonzero
nonidentity morphism lines $r$ of $T$, subject to the relations
$[rr']=[r]+[r']$ whenever $rr'\ne0$. If $s$ is a source of $T$ and
$T'=T\setminus\{s\}$, let $R$ be the representation-finite poset incidence
category obtained from $T'$ and $\rad T(s,-)$ as in
Lemma~\ref{lem:source-basis-extension}. Because the tautological weight
$r\mapsto[r]\in U(T')$ is additive on $R$, Lemma~\ref{lem:poset-weights}
provides elements $u_y\in U(T')$ for $y\in R$, normalized to $u_{y_D}=0$ at a
chosen base object $y_D$ in each connected component $D\in\pi_0(R)$, such
that $[b_{yz}]=u_z-u_y$ for every nonzero product $(sy)b_{yz}=sz$. Since
the relations in $U(T)$ involving the new generators $[sy]$ are precisely
$[sy]+[b_{yz}]=[sz]$, substituting $[sy]=e_D+u_y$ gives a canonical direct-sum
decomposition
\[
 U(T)\cong U(T')\oplus\Z^{\pi_0(R)},
\]
and the dual decomposition holds when adjoining a sink. Passing to the direct
limit along the one-object exhaustion $\mathcal L_1\subset\mathcal L_2\subset\cdots$
of Lemma~\ref{lem:cover-exhaustion} shows that $U(\mathcal L)$ is a free
abelian group.

Now consider the $G$-equivariant homomorphism
$\partial:U(\mathcal L)\to I_{\mathcal L}$ sending $[r:x\to y]$ to
$e_y-e_x$. Because $\mathcal L$ is connected, $\partial$ is surjective. We
claim that $\partial$ is also injective. Since $U(\mathcal L)$ is free
abelian and $k^*$ is infinite (so for any nonzero integer $m$ there exists
$a\in k^*$ with $a^m\ne1$), any nonzero element
$z\in\ker\partial$ would be detected by some group homomorphism
$\varphi:U(\mathcal L)\to k^*$; however, by
Lemma~\ref{lem:scalar-object-values}, every multiplicative weight
$\varphi\in\Hom(U(\mathcal L),k^*)=\operatorname{Scal}(\mathcal L)$ has the
form $\varphi([r:x\to y])=\zeta_y/\zeta_x$ and therefore factors through
$\partial$, forcing $\varphi(z)=1$. Thus $\partial:U(\mathcal L)\xrightarrow{\sim}I_{\mathcal L}$
is a $G$-equivariant isomorphism, which immediately implies (2).

Because the nonzero rays of $P$ and their nonzero products are the $G$-orbits
of the nonzero morphism lines of $\mathcal L$ and their nonzero products,
$W(P)$ is the group of $G$-coinvariants $U(\mathcal L)_G\cong(I_{\mathcal L})_G$.
Fix reference lifts $x_0,x_1,\ldots,x_{n-1}\in\operatorname{ob}\mathcal L$ of
the $n$ projective vertices of $A_0$. In $(I_{\mathcal L})_G$, the relation
$e_{gy}-e_x=(e_{gy}-e_{gx})+(e_{gx}-e_x)\equiv(e_y-e_x)+(e_{gx}-e_x)$ shows
that $\theta_g=[e_{gx}-e_x]$ is independent of $x$ and satisfies
$\theta_{gh}=\theta_g+\theta_h$. Thus $(I_{\mathcal L})_G$ is generated by
the classes $z_v=[e_{x_v}-e_{x_0}]$ ($1\leq v\leq n-1$) and $\theta_g$
($g\in G$), and the homomorphism $(I_{\mathcal L})_G\to\Z^{n-1}\oplus G_{\mathrm{ab}}$
induced by sending the object $gx_v$ to $(e_v,[g])$ (with $e_0=0$) provides
the inverse isomorphism, proving (1).
\end{proof}

\subsection{Integer weights on rays and meshes}

We use integer weights on $P$ and on $\C$ to refine the categorical grading
and separate first-order deformations attached to different periodic
components.

\begin{definition}
\label{def:integer-weights}
Let $A_0=kP$ ($P=\ray A_0$) and $\C=\gr_\JJ\ind(\modu A)$.
\begin{enumerate}[label={\upshape(\arabic*)}]
\item An \emph{integer ray weight} is a function $h:P\setminus\{0\}\to\Z$
that vanishes on identity rays and satisfies $h(rs)=h(r)+h(s)$ whenever
$rs\ne0$. We write $H_{\Z}=\Hom_{\Z}(W(P),\Z)$ for the group of integer ray
weights, set $L=W(P)$ (called the \emph{weight lattice} of $A_0$), and write
$\ell(r)=[r]\in L$ for the weight class of a nonzero ray $r\in P$.
\item An \emph{integer mesh weight} on $\C$ is a pair $w=(w,\nu)$ assigning an
integer $w(a)\in\Z$ to each arrow $a$ of $\Gamma_A$ and an integer
$\nu_v\in\Z$ to each nonprojective vertex $v$ of $\Gamma_A$ such that
\[
                   w(\sigma a)+w(a)=\nu_v\qquad\text{for every arrow }a:u\to v.
\]
We write $H^{\mathrm{mesh}}_{\Z}$ for the abelian group of integer mesh
weights on $\C$.
\end{enumerate}
\end{definition}

Note that the categorical degree function $\deg_\JJ$ of Lemma~\ref{lem:ray-degree}
is an element of $H_{\Z}$. Because $L=W(P)$ is free abelian of finite rank by
Proposition~\ref{prop:integral-ray-weights}(1), biduality identifies $L$
with $\Hom_{\Z}(H_{\Z},\Z)$ via $\chi(h)=h(\chi)$ for $\chi\in L$ and
$h\in H_{\Z}$. The algebra $A_0=kP$ is naturally $L$-graded by placing each
ray $r\in P$ in degree $\ell(r)\in L$; a Hochschild $2$-cochain sending a
pair of rays $(r,s)$ to a multiple of $u$ shifts $L$-degree by
$\ell(u)-\ell(r)-\ell(s)\in L$.

\begin{lemma}
\label{lem:integer-weight-lifting}
For $A_0=kP$ and $\C=\gr_\JJ\ind(\modu A)$, assigning to each path in $\C$ the
sum of the weights of its arrows under a mesh weight $w\in H^{\mathrm{mesh}}_{\Z}$
induces a well-defined grading on $\C$ and restricts on the projective
corners to an integer ray weight in $H_{\Z}$. This restriction fits into a
short exact sequence of abelian groups
\[
 0\longrightarrow
 \{(\zeta_X)\in\Z^{\operatorname{ob}\C}:\zeta_P=0
                         \text{ for all projective }P\}
 \longrightarrow H^{\mathrm{mesh}}_{\Z}\longrightarrow H_{\Z}
 \longrightarrow0,
\]
where the left-hand map sends a tuple $(\zeta_X)$ to the coboundary mesh
weight $w(a)=\zeta_{\operatorname{target}(a)}-\zeta_{\operatorname{source}(a)}$
and $\nu_v=\zeta_v-\zeta_{\tau v}$.
\end{lemma}
\begin{proof}
For any $w\in H^{\mathrm{mesh}}_{\Z}$, every term $(\sigma a)a$ of the mesh
relation at a nonprojective vertex $v$ has total weight $\nu_v$, so the mesh
ideal of $\C$ is homogeneous with respect to path weight. Because each
nonzero homogeneous component $A_0(x,y)_d=\C_d(P_y,P_x)$ is one-dimensional
by Lemma~\ref{lem:ray-degree}, all paths representing a given ray $r\in P$
have the same $w$-weight, and Proposition~\ref{prop:projective-ray-product}(2)
shows that this weight is additive on nonzero ray products. Thus restriction
defines a homomorphism $H^{\mathrm{mesh}}_{\Z}\to H_{\Z}$.

To see that restriction is surjective, pull any $h\in H_{\Z}$ back to an
additive weight on the category $\mathcal L$ of lifted projectives. By
Proposition~\ref{prop:integral-ray-weights}(2), there exist integers $b_x$
for $x\in\operatorname{ob}\mathcal L$ such that $h(r)=b_y-b_x$ for every
nonzero morphism $r\in\mathcal L(x,y)$, and by $G$-invariance the difference
$\gamma(g)=b_x-b_{gx}$ is independent of $x\in\operatorname{ob}\mathcal L$
and homomorphic in $g\in G$. On the lifted projective vertices of
$\widetilde\C$, set $\zeta_{P_x}=-b_x$ (the sign change reflecting that
$\mathcal L$ is the opposite category), and extend $\zeta$ to all vertices of
$\widetilde\C$ by choosing arbitrary integers on one representative of each
$G$-orbit of nonprojective vertices and setting
$\zeta_{g\widetilde X}=\zeta_{\widetilde X}+\gamma(g)$. Then the differences
of $\zeta$ along lifted arrows and translation edges are $G$-invariant and
therefore descend to an integer mesh weight on $\C$ whose restriction to $P$
is $\zeta_{P_x}-\zeta_{P_y}=b_y-b_x=h(r)$.

Finally, if $w\in H^{\mathrm{mesh}}_{\Z}$ restricts to zero on $P$, pull $w$
back to a closed integer edge weight on the simply connected complex
$\widetilde K$ to write $w$ as the difference of vertex integers
$\zeta_{\widetilde X}$. Because $w$ vanishes on all morphisms between lifted
projectives and $\mathcal L$ is connected, $\zeta$ is constant on all lifted
projective vertices; subtracting this constant makes $\zeta_{P_x}=0$ for all
$x$. Since $\zeta_{g\widetilde X}-\zeta_{\widetilde X}$ is independent of
$\widetilde X$ and vanishes on lifted projectives, $\zeta$ is $G$-invariant
and descends to a vertex assignment $(\zeta_X)$ on $\C$ vanishing at the
projective vertices. Because every connected component of $\Gamma_A$ contains
at least one projective vertex, this coboundary map is injective.
\end{proof}

For each residual component $i\in I$ and its representative path
$p_i:\tau v_i\to v_i$ of length $d_i=2s_i+2$ from
Proposition~\ref{prop:residual-correction-spaces}(2), define the linear
functional $\chi_i$ on $H^{\mathrm{mesh}}_{\Z}$ by
\begin{equation}
 \chi_i(w)=w(p_i)-\nu_{v_i}.
 \label{eq:parameter-character}
\end{equation}
Because $p_i$ and every term of the mesh relation at $v_i$ share the same
endpoints $\tau v_i$ and $v_i$, any vertex coboundary $(\zeta_X)$ shifts
$w(p_i)$ and $\nu_{v_i}$ by the same amount $\zeta_{v_i}-\zeta_{\tau v_i}$.
By Lemma~\ref{lem:integer-weight-lifting}, $\chi_i(w)$ therefore depends only
on the restricted ray weight in $H_{\Z}$ and defines an element
$\chi_i\in\Hom_{\Z}(H_{\Z},\Z)=L$, which we call the \emph{residual
character} of the $i$th residual component.

More generally, each mesh weight $w\in H^{\mathrm{mesh}}_{\Z}$ grades
$V_{m+1}(\C)$ by assigning weight $w(p)-\nu_v$ to a path $p:\tau v\to v$ in
row $v$, and grades $U_m(\C)$ by assigning weight $w(p)-w(a)$ to a path $p$
parallel to an arrow $a$. The mesh derivative $D_m$ of
\eqref{eq:mesh-derivative} preserves this grading, so
$\operatorname{coker}D_m(\C)$ becomes an $L$-graded space in which the line
$k[p_i]$ is homogeneous of $L$-degree $\chi_i$ (and of positive internal
degree $\chi_i(\deg_\JJ)=d_i-2=2s_i>0$).

\begin{proposition}
\label{prop:primitive-parameter-characters}
Assume $\charac k=2$. In the weight lattice $L=W(P)$ of
Definition~\ref{def:integer-weights}, the residual characters
$\{\chi_i\}_{i\in I}$ defined by \eqref{eq:parameter-character} extend to a
$\Z$-basis of $L$.
\end{proposition}
\begin{proof}
For each $i\in I$, identify the $i$th residual component with
$\Z\Delta/\langle\tau^{s_i}\rho\rangle$ as in
Proposition~\ref{prop:residual-correction-spaces}(1), with $\Delta$ oriented
away from the root fixed by $\rho$. On the vertices $(x,j)$ of $\Z\Delta$,
set
\[
                   F_i(x,j)=\left\lfloor j/s_i\right\rfloor,
\]
and assign to every arrow and translation edge of $\Z\Delta$ the difference
of $F_i$ between its target and source. Because the deck generator
$(x,j)\mapsto(\rho x,j-s_i)$ subtracts $1$ from $F_i$ at every vertex, these
edge differences are invariant under $\langle\tau^{s_i}\rho\rangle$ and
descend to a closed integer edge weight on the $i$th residual component.
Assigning weight zero to all other periodic components of $\Gamma_A$ and
extending to a closed edge weight on the entire mesh complex $K$ via
Lemma~\ref{lem:periodic-weight-extension} produces an integer mesh weight
$w_i\in H^{\mathrm{mesh}}_{\Z}$ on $\C$; let $h_i\in H_{\Z}=\Hom_{\Z}(L,\Z)$
be its restriction to $P$.

Lift the path $p_i:\tau v_i\to v_i$ of length $2s_i+2$ to $\Z\Delta$ starting
at a vertex $(x,r)$ over $\tau v_i$. Its endpoint has the form
$(\rho^m x,r+1-ms_i)$ for some $m\in\Z$ and lies at the same distance from
the root of $\Delta$ as $x$. Since outward arrows in $\Z\Delta$ preserve the
second coordinate while inward arrows increase it by $1$, the path $p_i$ of
length $2s_i+2$ traverses an equal number $s_i+1$ of outward and inward
arrows and therefore ends at second coordinate $r+s_i+1$. Subtracting the
weight of the translation edge $t_{v_i}$ from $(x,r)$ to $(x,r+1)$ gives
\[
 h_i(\chi_i)=\chi_i(w_i)
 =\left\lfloor\frac{r+s_i+1}{s_i}\right\rfloor
       -\left\lfloor\frac{r+1}{s_i}\right\rfloor=1,
\]
whereas $h_i(\chi_j)=\chi_j(w_i)=0$ for $j\ne i$ because $w_i$ vanishes on
the $j$th periodic component. Consequently, the homomorphism
$(h_i)_{i\in I}:L\to\Z^I$ is a split surjection with section sending the
$i$th standard basis vector to $\chi_i$, so
$L=\bigoplus_{i\in I}\Z\chi_i\oplus\ker(h_i)_{i\in I}$. Since $\ker(h_i)_{i\in I}$
is a subgroup of the finitely generated free abelian group $L$, it is free
abelian, and any $\Z$-basis of it completes $\{\chi_i\}_{i\in I}$ to a
$\Z$-basis of $L$.
\end{proof}

\subsection{Hochschild cohomology and first-order mesh deformations}
\label{subsec:hochschild-restriction}

We now relate the Hochschild cohomology of $A_0$ to the mesh derivative
cokernel by lifting the degree-normalized indecomposable $A_0$-modules to
first-order deformations of $A_0$ over $k[\varepsilon]/(\varepsilon^2)$.
Because the degree-zero vertex subalgebra $S=\prod_x ke_x$ is separable over
$k$, the inclusion of the normalized $S$-relative cochain complex
$C^\bullet_S(A_0,A_0)$ of \eqref{eq:normalized-relative-cochains} into the
ordinary Hochschild cochain complex is a quasi-isomorphism
\cite[Theorem~1.2]{GerstenhaberSchack}. Both the internal categorical degree
and the $L$-grading on $A_0=kP$ induce gradings on $\HH^2(A_0)$; we write
$\HH^2(A_0)_q$ for the component of internal degree $q$,
$\HH^2(A_0)_{>0}=\bigoplus_{q>0}\HH^2(A_0)_q$, and $\HH^2(A_0)_\chi$ for the
component of $L$-weight $\chi\in L$. (In
Proposition~\ref{prop:full-hochschild-comparison} we will show that the
injection $\Psi$ below is in fact an isomorphism and $I_0=I$, though only
injectivity is needed to prove Theorem~A.)

\begin{proposition}
\label{prop:hochschild-mesh-restriction}
Let $A_0$ be the standard form of Construction~\ref{con:projective-algebra}
and $\C$ its graded mesh category. For every internal degree $q>0$, lifting
degree-normalized modules to first-order deformations induces an $L$-graded
injection
\[
 \Psi:\HH^2(A_0)_q\hookrightarrow\operatorname{coker}D_{q+1}(\C).
\]
In particular, when $\charac k=2$, there is a subset $I_0\subseteq I$ of the
residual components of Proposition~\ref{prop:residual-correction-spaces}
such that
\[
 \HH^2(A_0)_{>0}=\bigoplus_{i\in I_0}\HH^2(A_0)_{\chi_i},
 \qquad \dim_k\HH^2(A_0)_{\chi_i}=1\quad(i\in I_0).
\]
\end{proposition}
\begin{proof}
Represent a cohomology class in $\HH^2(A_0)_q$ by a normalized $S$-relative
$2$-cocycle $\mu\in C^2_S(A_0,A_0)_q$, and form the first-order deformation
$A_{0,\mu}=(A_0\otimes_k k[\varepsilon]/(\varepsilon^2),m+\varepsilon\mu)$
with $\deg\varepsilon=-q$. By Lemma~\ref{lem:positive-ext-vanishing}(2) and
Lemma~\ref{lem:square-zero-selected-lifts}, the degree-normalized
indecomposable $A_0$-modules $M_1,\ldots,M_h$ lift to graded
$A_{0,\mu}$-modules $M_i^\mu$ (with $M_i^\mu=e_x A_{0,\mu}$ at the projective
vertices), and each arrow $a$ of $\Gamma_A$ lifts to a degree-one morphism
$\tilde a$ in the Hom category $\C_\mu$ of the lifted modules. Because the
mesh relation $r_v=\sum_{a:u\to v}(\sigma a)a$ vanishes in $\C$, its evaluation
on the lifted arrows $\tilde a$ in $\C_\mu$ vanishes modulo $\varepsilon$; by
the freeness of the lifted Hom spaces over $k[\varepsilon]/(\varepsilon^2)$
(Lemma~\ref{lem:square-zero-selected-lifts}(2)), this mesh error has the form
$\varepsilon z_v$ for a unique element $z_v\in\C_{q+2}(\tau v,v)$ (since
$\deg\varepsilon=-q$ and $\deg r_v=2$). Modifying the arrow lifts by
$\tilde a\mapsto\tilde a+\varepsilon f_a$ with $f_a\in\C_{q+1}$ changes the
tuple $z=(z_v)_v$ by $D_{q+1}(f)$, while changing the module lifts $M_i^\mu$
by automorphisms reducing to the identity or replacing $\mu$ by a cohomologous
cocycle $\mu+dh$ (which corresponds to applying the algebra automorphism
$1+\varepsilon h$ of $A_{0,\mu}$) leaves the class $[z]\in\operatorname{coker}D_{q+1}(\C)$
unchanged. Thus $[\mu]\mapsto\Psi([\mu])=[z]$ is well-defined.

Because the first-order module-action and arrow-lifting equations in
Lemma~\ref{lem:square-zero-selected-lifts} are linear in $\mu$ and in the
first-order corrections, and any product of two $\varepsilon$-corrections
vanishes modulo $\varepsilon^2$, the map $\Psi$ is $k$-linear. To see that
$\Psi$ is injective, suppose $\Psi([\mu])=0$. Then we can adjust the arrow
lifts $\tilde a$ by $\varepsilon f_a$ so that every mesh relation $r_v(\tilde a)$
vanishes identically in $\C_\mu$. The assignment $a\mapsto\tilde a$ then
defines a graded $k[\varepsilon]/(\varepsilon^2)$-linear functor
$\C\otimes_k k[\varepsilon]/(\varepsilon^2)\to\C_\mu$ that reduces to the
identity modulo $\varepsilon$. By Lemma~\ref{lem:square-zero-selected-lifts}(2),
the morphism spaces of $\C_\mu$ are free over $k[\varepsilon]/(\varepsilon^2)$
with bases lifting those of $\C$, so by Nakayama's lemma over the local ring
$k[\varepsilon]/(\varepsilon^2)$ this functor is an isomorphism on every Hom
space. Restricting to the projective vertices $e_x A_{0,\mu}$ yields a
graded algebra isomorphism $A_0\otimes_k k[\varepsilon]/(\varepsilon^2)\xrightarrow{\sim}A_{0,\mu}$
reducing to the identity modulo $\varepsilon$, which means that $[\mu]=0$ in
$\HH^2(A_0)_q$.

Finally, choosing an $L$-homogeneous cocycle $\mu$ of weight $\chi\in L$ and
lifting the integer ray weights to mesh weights on $\C$ via
Lemma~\ref{lem:integer-weight-lifting}, Lemma~\ref{lem:square-zero-selected-lifts}(3)
shows that when $\varepsilon$ is assigned $L$-weight $-\chi$, the lifted
modules and arrows can be chosen $L$-homogeneous, so the resulting mesh
error $z_v$ has weight $\chi$ relative to row $v$. Thus $\Psi$ preserves
$L$-weights. Combining the injectivity of $\Psi$ with
Proposition~\ref{prop:residual-correction-spaces}(2) immediately gives the
stated decomposition of $\HH^2(A_0)_{>0}$ in characteristic two.
\end{proof}
